\documentclass[11pt]{article}
\usepackage[margin=0.8in]{geometry}
\usepackage{xcolor,parskip,array,enumitem,listings}
\usepackage[hidelinks]{hyperref}
\definecolor{codegray}{gray}{0.97}
\definecolor{huddleblue}{RGB}{214,234,255}
\definecolor{predictionyellow}{RGB}{255,242,200}
\lstset{basicstyle=\ttfamily\small,columns=fullflexible,frame=single,
  framerule=0.4pt,backgroundcolor=\color{codegray},breaklines=true,
  showstringspaces=false,keepspaces=true,aboveskip=6pt,belowskip=6pt}
\setlist[enumerate]{leftmargin=*,itemsep=4pt,topsep=4pt}
\newcommand{\response}[2]{\par\noindent\fbox{\begin{minipage}{0.96\linewidth}\textbf{#1}\par\vspace{#2}\end{minipage}}\par}
\newcommand{\callout}[2]{\noindent\colorbox{#1}{\begin{minipage}{0.96\linewidth}#2\end{minipage}}\par}
\begin{document}
\begin{center}
{\LARGE MA 214: Lab 1}\\[3pt]
{\large Working with R and Data: NYC Flights}
\end{center}
\textbf{Name:} \hrulefill\quad  \textbf{BU ID:} \hrulefill\par
\textbf{Partner's name:} \hrulefill \quad \textbf{BU ID:} \hrulefill

\textbf{Mission:} Use flight records to ask a question, choose a summary, make a plot, and support a claim. Work with a partner and explain your code to each other.
\textbf{Do not use GenAI tools for this lab.} Practice your own coding and statistical reasoning.

\section*{Getting Started: Packages, Data, and Help}
Open \texttt{lab-starter.R}. We will use the full \texttt{flights} data set in \texttt{nycflights13}.

Run the installation commands in the \textbf{Console once}, if needed:
\begin{lstlisting}
install.packages("nycflights13")
install.packages(c("dplyr", "ggplot2"))
\end{lstlisting}
Then run these lines at the start of each new R session:
\begin{lstlisting}
library(nycflights13)
library(dplyr)
library(ggplot2)
data(flights)
?flights
\end{lstlisting}
\textbf{What do these commands do?}
\begin{center}
\renewcommand{\arraystretch}{1.25}
\begin{tabular}{p{0.35\linewidth}p{0.57\linewidth}}
\textbf{Command} & \textbf{Purpose} \\
\texttt{install.packages("name")} & Downloads and installs a package on your computer. \\
\texttt{library(name)} & Attaches an installed package for this session. \\
\texttt{data(flights)} & Loads the packaged data set for use in R. \\
\texttt{?flights} & Opens documentation: context, variables, and units. \\
\texttt{head(flights)} & Shows the first six rows. \\
\texttt{dim(flights)} & Returns the number of rows and columns. \\
\texttt{names(flights)} & Lists the column names. \\
\texttt{glimpse(flights)} & Shows column types and example values. \\
\end{tabular}
\end{center}
\callout{huddleblue}{\textbf{Huddle:} Why do we install a package once but load it again in a new session? What information can the help page give you that the first six rows cannot?}
\response{Warm-up: Introduce yourself. Which city did you visit recently? If you flew, how much time did you spend at the airport?}{0.35in}

\newpage
\section*{1. Meet the Data and Read a Pipeline}
Run \texttt{head()}, \texttt{dim()}, \texttt{names()}, and \texttt{glimpse()} on \texttt{flights}. Use \texttt{?flights} to interpret the output.
\begin{enumerate}
\item How many rows and columns are there? What does one row represent? Which year and origin airports are covered?
\item For each variable below, identify whether it is categorical or quantitative and explain its meaning (including units). Is every column stored as a number a quantitative measurement?
\end{enumerate}
\begin{center}
\renewcommand{\arraystretch}{1.65}
\begin{tabular}{|p{0.20\linewidth}|p{0.28\linewidth}|p{0.41\linewidth}|}
\hline
\textbf{Variable} & \textbf{Type} & \textbf{Meaning / units} \\
\hline
\texttt{origin}, \texttt{dest} & & \\
\hline
\texttt{carrier}, \texttt{flight} & & \\
\hline
\texttt{hour} & & \\
\hline
\texttt{arr\_delay} & & \\
\hline
\texttt{air\_time}, \texttt{distance} & & \\
\hline
\end{tabular}
\end{center}
\textbf{Reading R:} \texttt{<-} saves a result with a name; \texttt{==} tests equality; \texttt{\%>\%} passes a result into the next function. Read a pipeline as ``take this, then do that.''
\begin{lstlisting}
airport_counts <- flights %>%
  count(origin, sort = TRUE)
airport_counts
\end{lstlisting}
\begin{center}
\renewcommand{\arraystretch}{1.2}
\begin{tabular}{p{0.38\linewidth}p{0.54\linewidth}}
\textbf{Function} & \textbf{Purpose} \\
\texttt{filter(origin == "JFK")} & Keeps rows matching a condition. \\
\texttt{select(origin, hour)} & Keeps selected columns. \\
\texttt{count(origin)} & Counts rows in each origin group. \\
\texttt{arrange(desc(n))} & Orders rows from largest to smallest \texttt{n}. \\
\texttt{group\_by(origin)} & Sets the groups for the next summary. \\
\texttt{summarize(...)} & Calculates summaries within each group. \\
\texttt{n()}, \texttt{mean()}, \texttt{median()} & Group size, average, and middle value. \\
\end{tabular}
\end{center}
\response{Which airport has the most records? Then select three variables and filter to one airport in your script.}{0.4in}
\textbf{Watch the units:} \texttt{hour} is the \emph{scheduled} departure hour (local time). \texttt{dep\_time} uses HHMM, so 830 means 8:30, not 830 minutes. Negative delay means early; \texttt{NA} means missing, not zero.

\newpage
\section*{2. When Is an Airport Busiest?}
\callout{predictionyellow}{\textbf{Predict first:} Choose JFK, LGA, or EWR. Which scheduled departure hour do you expect to have the most flights, and why?}
\response{Chosen airport, predicted busiest hour, and reason:}{0.3in}
Here, \textbf{busiest hour} means the hour of day with the most scheduled departures in the data across all of 2013. We are counting records, not passengers or simultaneous aircraft.
\begin{lstlisting}
chosen_airport <- "JFK"  # Change this to your airport.
hour_counts <- flights %>%
  filter(origin == chosen_airport) %>%
  count(hour, name = "n_flights") %>%
  arrange(desc(n_flights))
hour_counts
\end{lstlisting}
\begin{enumerate}
\item Report the busiest hour as a clock interval (for example, 9 means 9:00--9:59) and its flight count. Check for ties. Did the result support your prediction?
\item Repeat for another airport. Do the peak hours match? What information is lost when we combine all days of the year?
\end{enumerate}
\response{Evidence: Compare your two airports using hours and counts.}{0.5in}
\textbf{Make a plot.} \texttt{ggplot()} starts a plot; \texttt{aes()} maps columns to axes; \texttt{geom\_col()} draws bars using the counts we already calculated.
\begin{lstlisting}
hour_plot <- ggplot(hour_counts,
                    aes(x = hour, y = n_flights)) +
  geom_col() +
  scale_x_continuous(breaks = seq(0, 23, by = 3)) +
  labs(title = paste("Scheduled departures from", chosen_airport),
       x = "Scheduled departure hour (local time)",
       y = "Flight count in 2013")
hour_plot
# Optional: save to your current working directory.
# ggsave("lab1-hour-counts.png", hour_plot, width = 7, height = 4)
\end{lstlisting}
\callout{huddleblue}{\textbf{Huddle:} Why does a histogram of \texttt{dep\_time} answer a different question? Why does a larger flight count not tell us the number of passengers?}

\newpage
\begingroup\setlength{\parskip}{4pt}
\section*{3. How Long Is the Delay?}
Which origin airport has the greatest mean arrival delay? Report the mean for EWR, JFK, and LGA.

\textit{Delays are in minutes; negative values mean early arrivals. Enter and run this code in your script. Missing values are excluded.}

\begin{lstlisting}
airport_delays <- flights %>%
  group_by(origin) %>%
  summarize(
    n_observed = sum(!is.na(arr_delay)),
    mean_delay = mean(arr_delay, na.rm = TRUE),
    .groups = "drop"
  ) %>%
  arrange(desc(mean_delay))
airport_delays
\end{lstlisting}

\response{Your answer:}{0.75in}

\section*{4. Can a Late Departure Still Arrive on Time?}
Of flights departing more than 30 minutes late, what percentage arrive on time or early?

\textit{Use \texttt{dep\_delay > 30} and \texttt{arr\_delay <= 0}. Exclude missing arrival delays; report the count and percentage.}

\begin{lstlisting}
recovery <- flights %>%
  filter(dep_delay > 30, !is.na(arr_delay)) %>%
  summarize(
    n_observed = n(),
    n_recovered = sum(arr_delay <= 0),
    pct_recovered = 100 * n_recovered / n_observed
  )
recovery
\end{lstlisting}

\response{Your answer:}{0.7in}

\newpage
\section*{5. Late Leaving, Late Arriving?}
Make a scatterplot of departure delay and arrival delay. What relationship do you see?

\textit{Use \texttt{dep\_delay} on the x-axis and \texttt{arr\_delay} on the y-axis (both in minutes). Each point is a flight. Describe the pattern in 1--2 sentences.}

\response{Sketch your plot and describe the relationship:}{1.5in}

\section*{6. What Else Would You Like to Find Out?}
Write one question you could investigate with these data and name the variables you would use.

\response{Your answer:}{0.7in}

\textbf{Before you leave:} Submit the completed activity worksheet.

{\footnotesize Data documentation: \url{https://nycflights13.tidyverse.org/reference/flights.html}.}
\par\endgroup
\end{document}
